\section{Historické metody}\label{sec:krypto-historicke}

Nejstarší šifry pracovaly přímo se znaky textu. Zprávu bylo možné šifrovat ručně bez počítače, jejich jednoduchost je však zároveň slabinou. Na historických metodách dobře uvidíme, proč velký počet klíčů sám o~sobě ještě nezaručuje bezpečnost.

\subsection{Substituční šifra s~posunem}

Mějme abecedu
\[
    \Sigma=\set{a_0,a_1,\ldots,a_{q-1}}
\]
a~tajný posun $k$, kde $0<k<q$. Každý znak nahradíme znakem posunutým o~$k$ míst:
\[
    a_i\longmapsto a_{(i+k)\bmod q}.
\]
Operace modulo $q$ způsobí, že se po dosažení konce abecedy vrátíme na její začátek. Dešifrování používá opačný posun:
\[
    a_j\longmapsto a_{(j-k)\bmod q}.
\]

\begin{figure}[ht]
    \centering
    \begin{tikzpicture}[
    font=\small,
    cell/.style={draw=black!55, rounded corners=.6mm,
                 minimum width=8mm, minimum height=8mm, inner sep=0pt},
    arrow/.style={-{Latex}, draw=black!65, line width=.7pt}
]
    \foreach \x/\letter in {0/A,1/B,2/C,3/D,4/E,5/F,6/G,7/H}{
        \node[cell, fill=black!4] (u\x) at (\x,0) {\letter};
    }
    \foreach \x/\letter in {0/D,1/E,2/F,3/G,4/H,5/A,6/B,7/C}{
        \node[cell, fill=green!15] (d\x) at (\x,-1.4) {\letter};
    }
    \foreach \x in {0,...,7}{
        \draw[arrow] (u\x) -- (d\x);
    }
    \node[anchor=east] at (-.65,0) {původní znak};
    \node[anchor=east] at (-.65,-1.4) {šifrový znak};
    \node[font=\small\bfseries] at (3.5,.8) {posun \(k=3\)};
\end{tikzpicture}

    \caption{Substituční šifra s~posunem $k=3$ na zkrácené abecedě.}
    \label{fig:krypto-caesar}
\end{figure}

Pro standardní latinskou abecedu a~$k=3$ například dostaneme
\[
    \texttt{TAJNE}\longmapsto\texttt{WDMQH}.
\]
Jestliže záškodník nezná klíč, může jednoduše vyzkoušet všech $q-1$ nenulových posunů. Pro $q=26$ jde pouze o~25 možností. U~dostatečně dlouhé zprávy navíc snadno pozná, který výsledek připomíná smysluplný text.

\subsection{Obecná substituční šifra}

Místo jednoho posunu můžeme zvolit libovolnou permutaci abecedy
\[
    \mapping{\pi}{\Sigma}{\Sigma}.
\]
Permutace je bijekce, takže každý znak má právě jeden zašifrovaný obraz a~každý znak šifrové abecedy má právě jeden vzor. Šifrování nahradí každý znak $a$ znakem $\pi(a)$, zatímco dešifrování používá inverzní permutaci $\pi^{-1}$.

Pro abecedu o~$q$ znacích existuje $q!$ různých permutací. U~26 písmen je to přibližně $4\cdot 10^{26}$ klíčů, takže prosté vyzkoušení všech možností už není praktické. Obecná substituce však stále zachovává strukturu původního textu: stejné znaky se vždy mění na stejné znaky, opakovaná slova vytvářejí stejné vzory a~četnosti písmen se pouze přejmenují.

\begin{figure}[ht]
    \centering
    \input{06-kryptografie/images/ch06-substituce.tex}
    \caption{Obecná substituce a~zachování četností znaků.}
    \label{fig:krypto-substituce}
\end{figure}

Uvažujme například klíč, jehož prvních několik přiřazení je
\[
    A\mapsto Q,\quad B\mapsto W,\quad C\mapsto E,\quad
    D\mapsto R,\quad E\mapsto T,\quad\ldots
\]
Pro permutaci použitou na obrázku~\ref{fig:krypto-substituce} dostáváme
\[
    \texttt{TAJEMSTVI}\longmapsto\texttt{ZQPTDLZCO}.
\]

\warningbox{Velký počet klíčů neznamená automaticky bezpečnou šifru. Obecná substituce má obrovský prostor klíčů, ale frekvenční analýza využívá vlastnosti jazyka a~nemusí jednotlivé klíče zkoušet jeden po druhém.}

\paragraph{Frekvenční analýza.}

V~delším českém textu se některá písmena objevují výrazně častěji než jiná. Jestliže se v~šifrovém textu mimořádně často vyskytuje znak \texttt{Q}, může odpovídat některému běžnému písmenu otevřeného textu. Záškodník dále sleduje dvojice a~trojice znaků, délky slov nebo opakované vzory. Každý takový poznatek zmenšuje množinu možných permutací.

Historické metody nám proto ukazují dvě důležité zásady. Bezpečnost nelze odvozovat jen z~toho, že klíčů je mnoho, a~šifra by neměla zbytečně zachovávat statistickou strukturu otevřeného textu. Moderní šifry se snaží, aby bez znalosti klíče vypadal šifrový text jako náhodná data.
